\documentclass[a4paper,12pt]{article}\usepackage[]{graphicx}\usepackage[]{color}
%% maxwidth is the original width if it is less than linewidth
%% otherwise use linewidth (to make sure the graphics do not exceed the margin)
\makeatletter
\def\maxwidth{ %
  \ifdim\Gin@nat@width>\linewidth
    \linewidth
  \else
    \Gin@nat@width
  \fi
}
\makeatother

\definecolor{fgcolor}{rgb}{0.345, 0.345, 0.345}
\newcommand{\hlnum}[1]{\textcolor[rgb]{0.686,0.059,0.569}{#1}}%
\newcommand{\hlstr}[1]{\textcolor[rgb]{0.192,0.494,0.8}{#1}}%
\newcommand{\hlcom}[1]{\textcolor[rgb]{0.678,0.584,0.686}{\textit{#1}}}%
\newcommand{\hlopt}[1]{\textcolor[rgb]{0,0,0}{#1}}%
\newcommand{\hlstd}[1]{\textcolor[rgb]{0.345,0.345,0.345}{#1}}%
\newcommand{\hlkwa}[1]{\textcolor[rgb]{0.161,0.373,0.58}{\textbf{#1}}}%
\newcommand{\hlkwb}[1]{\textcolor[rgb]{0.69,0.353,0.396}{#1}}%
\newcommand{\hlkwc}[1]{\textcolor[rgb]{0.333,0.667,0.333}{#1}}%
\newcommand{\hlkwd}[1]{\textcolor[rgb]{0.737,0.353,0.396}{\textbf{#1}}}%
\let\hlipl\hlkwb

\usepackage{framed}
\makeatletter
\newenvironment{kframe}{%
 \def\at@end@of@kframe{}%
 \ifinner\ifhmode%
  \def\at@end@of@kframe{\end{minipage}}%
  \begin{minipage}{\columnwidth}%
 \fi\fi%
 \def\FrameCommand##1{\hskip\@totalleftmargin \hskip-\fboxsep
 \colorbox{shadecolor}{##1}\hskip-\fboxsep
     % There is no \\@totalrightmargin, so:
     \hskip-\linewidth \hskip-\@totalleftmargin \hskip\columnwidth}%
 \MakeFramed {\advance\hsize-\width
   \@totalleftmargin\z@ \linewidth\hsize
   \@setminipage}}%
 {\par\unskip\endMakeFramed%
 \at@end@of@kframe}
\makeatother

\definecolor{shadecolor}{rgb}{.97, .97, .97}
\definecolor{messagecolor}{rgb}{0, 0, 0}
\definecolor{warningcolor}{rgb}{1, 0, 1}
\definecolor{errorcolor}{rgb}{1, 0, 0}
\newenvironment{knitrout}{}{} % an empty environment to be redefined in TeX

\usepackage{alltt}
\usepackage[applemac]{inputenc}
\usepackage[OT1,T1]{fontenc}
% \usepackage{times}
\usepackage{setspace}
\usepackage{lscape}
\usepackage{amsmath,amssymb}
\usepackage{fancyhdr}
\usepackage{fancyvrb,moreverb,boxedminipage}
\usepackage{multirow}
\usepackage{float}
\usepackage{hyperref}
\usepackage{listings}
\usepackage{verbatim}
\usepackage{enumitem}

\setlength{\textwidth}{160mm}
\setlength{\textheight}{230mm}
\setlength{\oddsidemargin}{-5mm}
\setlength{\topmargin}{-10mm}

\textwidth=6in
\textheight=8in
\topmargin=0.0in
\oddsidemargin=0in 
\baselineskip=18pt
\headheight=2cm
\setcounter{MaxMatrixCols}{10}


\newcommand{\ns}{\!\!\!}
\newcommand{\e}{\varepsilon}
\newcommand{\be}{\beta}
\newcommand{\s}{\sigma}
\newcommand{\vv}{\vspace{.5cm}}
\newcommand{\beps}{\boldsymbol{\epsilon}}
\newcommand{\sumin}{\sum_{i=1}^n}
\newcommand{\bg}[1]{\mbox{\boldmath{$#1$}}}
    \DeclareMathOperator{\plim}{plim}
  \DefineVerbatimEnvironment{code}{Verbatim}{fontsize=\small}
  \DefineVerbatimEnvironment{example}{Verbatim}{fontsize=\small}
  \newenvironment{answer}{\color{blue}}
  
  \def\by{\mathbf{y}}
  \def\bx{\mathbf{x}}
  \def\ba{\mathbf{a}}
  \def\bb{\mathbf{b}}
  \def\be{\mathbf{e}}
  \def\bu{\mathbf{u}}
  \def\bv{\mathbf{v}}
  \def\bz{\mathbf{z}}
  
  %===========
    %Greek letter vectors
  %===========
    \def\bbeta{\mbox{\boldmath $\beta$}}
  \def\bgamma{\mbox{\boldmath $\gamma$}}
  \def\bvareps{\mbox{\boldmath $\varepsilon$}}
  \def\bleta{\mbox{\boldmath $\eta$}}
  \def\bmu{\mbox{\boldmath $\mu$}}
  \def\btheta{\mbox{\boldmath $\theta$}}
  \def\bsigma{\mbox{\boldmath $\sigma$}}
  \def\blambgam{\mbox{\boldmath $\lambda_{\gamma}$}}
  \def\blambbet{\mbox{\boldmath $\lambda_{\beta}$}}
  \def\blambda{\mbox{\boldmath $\lambda$}}
  
  %===========
    %Matrices
  %===========
    \def\bX{\mathbf{X}}
  \def\bD{\mathbf{D}(\bsigma)}
  \def\bV{\mathbf{V}(\bsigma)}
  \def\bVinv{\mathbf{V}^{-1}(\bsigma)}
  \def\bR{\mathbf{R}(\bsigma)}
  \def\bA{\mathbf{A}}
  \def\bB{\mathbf{B}}
  \def\bZ{\mathbf{Z}}
  \def\bIn{\mathbf{I_{N}}}
  \def\bIm{\mathbf{I_{m}}}
  \def\bIni{\mathbf{I_{n}}}
  \def\SSW{\mathrm{SSW}}
  \def\SSB{\mathrm{SSB}}
  \def\Normal{\mathcal{N}\!}
  
  %===========
    %Parentheses, etc.
  %===========
    \def\op{\left(}
  \def\cp{\right)}
  \def\oc{\left[}
  \def\cc{\right]}
  \def\ok{\left\{}
    \def\ck{\right\}}
  
  %===========
    %Statistical functions and real numbers.
  %===========
    \def\exE{\mathbb{E}}
  \def\Real{\mathbb{R}}
  \def\Var{\mathbb{V}ar}
  \def\Proba{\mathbb{P}}
  \def\Cov{\mathbb{C}ov}
  
  %===========
    %Listings and hyperlink settings
  %===========
    \hypersetup{colorlinks=true,linkcolor=blue,citecolor=blue, urlcolor=blue}
  \urlstyle{tt}
  
  %===========
    %Special Environments
  %===========
    \newenvironment{exercices}{\newcounter{moncompteur}\begin{list}
      {\bfseries Ex. \arabic{moncompteur}}{\usecounter{moncompteur}
        \setlength{\labelwidth}{2.6em}\setlength{\leftmargin}{3.1em}}}{\end{list}}
  \renewcommand{\labelenumi}{\bfseries\alph{enumi})}
  \newcommand{\splus}[3][0]{\texttt{>~#2}\ \ \ \ \textit{#3}\\[#1mm]}
      
      \newenvironment{exo}
    {\begin{center}\setlength{\textwidth}{140mm} \underline{Exercise:} \begin{minipage}[t]{120mm}}
    {\end{minipage}\setlength{\textwidth}{160mm}\end{center}}
    
    
    \usepackage{verbatim}
    
    \lstset{ %
      language=R, 
      basicstyle=\ttfamily,
      backgroundcolor=\color{lightgray},   % choose the background color; you must add \usepackage{color} or \usepackage{xcolor}
      xleftmargin= 10 pt,
      xrightmargin= 10 pt,
      escapeinside={\%*}{*)},          % if you want to add LaTeX within your code
        frame=none,                    % adds a frame around the code
        numbers=none,                    % where to put the line-numbers; possible values are (none, left, right)
        numbersep=5pt,                   % how far the line-numbers are from the code
        numberstyle=\tiny\color{mygray}, % the style that is used for the line-numbers
        showspaces=false,                % show spaces everywhere adding particular underscores; it overrides 'showstringspaces'
      }
      
      \pagestyle{fancy}
      
      
      \definecolor{mygreen}{rgb}{0,0.6,0}
      \definecolor{mygray}{rgb}{0.5,0.5,0.5}
      \definecolor{mymauve}{rgb}{0.58,0,0.82}
      \definecolor{lightgray}{gray}{0.95}
\IfFileExists{upquote.sty}{\usepackage{upquote}}{}
\begin{document}
      
      
      \lhead{University of Geneva\\ \textbf{GLAM}\\
      Pr. Eva Cantoni}
      \rhead{GSEM\\ Spring 2021 \\ \textbf{Homework 10}}
      
      \begin{center}
      \bigskip
      
      \bigskip
      
      \bigskip
      
      \bigskip
      
      \bigskip
      
      \bigskip
      
      \fbox{{\Large Non-parametric regression II -- Handout }}
      
      \bigskip
      
      \bigskip
      
      \bigskip
      \end{center}
      
      %%%%%%%%%%%%%%%%%%%%%%%%%%%%%%%%%%%%%%%%%%%%%%%%%%%%%%%%%%%%%%%%%%%%%%%%%%%%
        
        
        \section{\underline{Exercise 2}}
      
      Let us consider a univariate nonparametric model. As shown during the lesson (\textit{c.f.} slide 271), the basis functions allows the linear representation:
        \begin{equation*}
      \by = \bX \bbeta + \mathbf{\epsilon}.
      \end{equation*}
      In order to control the wiggliness, the parameters $\bbeta$ are estimated \textit{via} minimization of a penalized cost function:
      \begin{equation*}
      \mathcal{C}\op\bbeta\cp = \| \by-\bX\bbeta\|_2^2 + \lambda\bbeta^T \mathbf{S}\bbeta \\.
      \end{equation*}
      In this exercise, we reparametrize the model in order to better understand some of its features. The so-called \textit{natural parametrization} is described in the course (slide 292).
      \begin{enumerate}
      \item Following the course, define $\mathbf{P} = \mathbf{U}^T\mathbf{R}$, $\bbeta' = \mathbf{P} \bbeta$ and $\bX' = \mathbf{Q}\mathbf{U}$ (allowing $\bX = \bX' \mathbf{P}$), and express the cost function $\mathcal{C}\op\bbeta'\cp$ under this new specification.
      
  \begin{answer}
  \begin{align*}
  \mathcal{C}\op \bbeta' \cp&= \| \by-\bX\bbeta\|_2^2 + \lambda\bbeta^T\mathbf{S}\bbeta = \| \by-\underbrace{\mathbf{Q}\mathbf{U}}_{\mathbf{\bX'}}\underbrace{\mathbf{U}^T\mathbf{R}\bbeta}_{\bbeta'}\|_2^2 + \lambda\bbeta^T\mathbf{R}^T\underbrace{\mathbf{R}^{-T}\mathbf{S}\mathbf{R}^{-1}}_{\mathbf{U}\mathbf{\Lambda}\mathbf{U}^{T}}\mathbf{R}\bbeta\\&= \| \by-\bX'\bbeta'\|_2^2 +  \lambda\underbrace{\bbeta^T\mathbf{R}^T\mathbf{U}}_{\bbeta'^T} \mathbf{\Lambda}\underbrace{\mathbf{U}^{T} \mathbf{R}\bbeta}_{\bbeta'}= \| \by-\bX'\bbeta'\|_2^2 + \lambda \bbeta'^T \mathbf{\Lambda}\bbeta'
\end{align*}
\end{answer}
\item Provide the estimator $\widehat{\bbeta'}$ of $\bbeta'$ minimizing $\mathcal{C}\op\bbeta'\cp$.
                                                                                                                                                     
\begin{answer}
\begin{align*}
\mathcal{C}\op\bbeta'\cp&= \| \by-\bX'\bbeta'\|_2^2 + \lambda \bbeta'^T \mathbf{\Lambda}\bbeta'= \op\by-\bX'\bbeta'\cp^T \op\by-\bX'\bbeta'\cp + \lambda\bbeta'^T \mathbf{\Lambda}\bbeta' \\& =\by^T\by - \bbeta'^T\bX'^T\by - \by^T\bX'\bbeta' +\bbeta'^T\bX'^T\bX'\bbeta'+\lambda\bbeta'^T\mathbf{\Lambda}\bbeta'\\&= \by^T\by - \bbeta'^T\bX'^T\by - \by^T\bX'\bbeta' +\bbeta'^T\oc \bX'^T\bX' + \lambda\mathbf{\Lambda}\cc\bbeta' \\\frac{\partial\mathcal{C}}{\partial \bbeta'} \op \bbeta' \cp&=  -2\bX'^T\by + 2\oc \bX'^T\bX' + \lambda\mathbf{\Lambda} \cc\bbeta'         
\end{align*}

Finally :
      
\begin{equation*}
\frac{\partial\mathcal{C}}{\partial \bbeta'} \op \widehat{\bbeta'} \cp = 0 \Leftrightarrow  -2\bX'^T\by + 2\oc \bX'^T\bX' + \lambda\mathbf{\Lambda} \cc\widehat{\bbeta'} = 0 \Leftrightarrow \widehat{\bbeta'} = \oc \bX'^T\bX' + \lambda\mathbf{\Lambda} \cc^{-1}\bX'^T\by
\end{equation*}
\end{answer}

\item Deduce the linear smoother $\mathbf{A}$ such that $\widehat{y} = \mathbf{A} y$. Compute $\mathrm{tr}(\mathbf{A})$ and find an interpretation for this quantity.

\begin{answer}
        Thus, this penalized regression smoother provides the following \emph{fitted values}: 
        \begin{equation*}
        \widehat{y} = \bX'\widehat{\bbeta'} 
        = \underbrace{\bX' \oc \mathbf{I} + \lambda\mathbf{\Lambda} \cc^{-1}\bX'^T}_{\mathbf{A}}\by 
        = \mathbf{A} y
        \end{equation*}
        The matrix $\mathbf{A}$ is thus a projection (or \emph{hat}) matrix with rank equal to its trace. Thus, the \emph{effective degrees of freedom} of the smoothing procedure can be computed via the computation of this trace. 
        \begin{align*}
        \textrm{tr}(\mathbf{A}) 
        &= \textrm{tr}\op \bX' \oc \mathbf{I} + \lambda\mathbf{\Lambda} \cc^{-1}\bX'^T\cp
        = \textrm{tr}\op \bX'^T\bX' \oc \mathbf{I} + \lambda\mathbf{\Lambda} \cc^{-1}\cp
        = \textrm{tr}\op \oc \mathbf{I} + \lambda\mathbf{\Lambda} \cc^{-1} \cp \\
        \textrm{tr}(\mathbf{A})  &= \sum_{k=0}^{K} \frac{1}{1+\lambda \Lambda_{k+1,k+1}} < K+1
        \end{align*}     
    This value is smaller than $K+1$, since the values in the denominator are larger than 1. 

    \end{answer}

\item Explain what happens to $\mathbf{A}$ as $\lambda$ increases.

\begin{answer}
    As $\lambda$ increases, we reduce the effective degrees of freedom in the smooth. When $\lambda$ is close to zero, the effective degrees of freedom are equal to $K+1$, meaning that all parameters have been used to represent the smooth. On the other hand, it is not excluded that this sum adds up to $1$, implying that a linear fit would be sufficient to model the impact of the covariate on the response. 
    
    The intepretation of the \emph{edf} extrapolates to the estimates obtained with Generalized Additive Models with Smoothing Parameters obtained by minimizing a criterion (UBRE, GCV, etc.). 
    \end{answer}

\end{enumerate}
    
\end{document}
    
